\documentclass{book}
% \documentclass{EngC}

%\usepackage[spanish]{babel}

\renewcommand{\chaptername}{Capítulo}
\renewcommand{\figurename}{Figura}

\makeatletter
\@ifclassloaded{EngC}{
    % Ensure standard size switches exist for packages expecting them
    \providecommand{\LARGE}{\Large}
    \providecommand{\Huge}{\huge}
}{
    % \usepackage[margin=1in]{geometry}
    \usepackage{fancyhdr}
    \pagestyle{fancy}
    \fancyhf{}
    \fancyhead[LE]{\thepage}
    \fancyhead[RE]{\ifnum\value{chapter}>0 \@chapapp\ \thechapter\fi}
    \fancyhead[LO]{\ifnum\value{chapter}>0 \S\thesection\fi}
    \fancyhead[RO]{\thepage}
}
\makeatother

% Font packages
\usepackage[utf8]{inputenc}
\usepackage{bbold}
\usepackage{pifont}

% Math packages
\usepackage{amsmath}
\usepackage{amssymb}

% Figures and colors
\usepackage{xcolor}
\usepackage{graphicx}
\usepackage{caption}
\usepackage{float}
\usepackage{subcaption}

% Diagrams and images
\usepackage{tikz} % General diagram software
\usepackage{tikz-cd} % Commutative diagrams
\usepackage{pgfplots} % PGF plots; used to plot functions
\usetikzlibrary{automata, arrows, angles, calc, intersections, math,
decorations.pathmorphing, decorations.markings}
\usetikzlibrary{positioning}
%% the following commands are needed for some matlab2tikz features
\usetikzlibrary{plotmarks}
\usetikzlibrary{arrows.meta}
\pgfplotsset{compat=1.16}
\usepgfplotslibrary{groupplots} % This lets you put multiple
% functions together in a single plot
\usepgfplotslibrary{patchplots}

% Tables
\usepackage{booktabs}
\usepackage{colortbl}
\usepackage{multirow}

% References and bibliography
\usepackage[
    backend=biber,
    style=alphabetic,
    sorting=nyt,
    doi=false,
    hyperref=true,
    isbn=false,
    sortcites,
    url=false,
    maxbibnames=99,
    maxalphanames=3,
    minalphanames=3,
    maxcitenames=2,
    mincitenames=1,
    backref=false,
    natbib=true,
]{biblatex}
\renewbibmacro{in:}{}
\AtEveryBibitem{%
    \clearlist{language}%
}
% Hack to accommodate texlab users
\providecommand{\toplevelprefix}{.}  % necessary for subfile
% compilation to work with bibliography
\makeatletter
\@ifundefined{thiscommandisneverdefined}{
    \addbibresource{\toplevelprefix/book-references.bib}
}{
    \addbibresource{book-references.bib}
}
\makeatother

\usepackage{hyperref}
\hypersetup{
    colorlinks=true,
    citecolor={blue!80!black},
}

\usepackage{cleveref}

\crefname{appendix}{apéndice}{apéndices}
\Crefname{appendix}{Apéndice}{Apéndices}
\crefname{assumption}{supuesto}{supuestos}
\Crefname{assumption}{Supuesto}{Supuestos}
\crefname{example}{ejemplo}{ejemplos}
\Crefname{example}{Ejemplo}{Ejemplos}
\crefname{exercise}{ejercicio}{ejercicios}
\Crefname{exercise}{Ejercicio}{Ejercicios}
\crefname{equation}{ecuación}{ecuaciones}
\Crefname{equation}{Ecuación}{Ecuaciones}
\crefname{figure}{figura}{figuras}
\Crefname{figure}{Figura}{Figuras}
\crefname{chapter}{capítulo}{capítulos}
\Crefname{chapter}{Capítulo}{Capítulos}
\crefname{section}{sección}{secciones}
\Crefname{section}{Sección}{Secciones}
\crefname{thm}{teorema}{teoremas}
\Crefname{thm}{Teorema}{Teoremas}
\crefname{remark}{observación}{observaciones}
\Crefname{remark}{Observación}{Observaciones}

\hyphenation{line-break line-breaks docu-ment triangle cambridge amsthdoc
cambridgemods baseline-skip author authors cambridgestyle en-vir-on-ment polar}

\setcounter{tocdepth}{2}% list subsections in the table of contents
% Algorithm packages -- should go after hyperref
\usepackage{algorithm}
% \usepackage[spaceRequire=false]{algpseudocodex}
\usepackage{algpseudocode}

% Math formatting
\usepackage{mathtools}
\usepackage{grffile}

% Text formatting
\usepackage[shortlabels]{enumitem}
\setlist[enumerate]{label=\arabic*., leftmargin=*}
\usepackage{fancyvrb}
\usepackage{soul}
\VerbatimFootnotes
\allowdisplaybreaks

% Custom packages for the book
\usepackage{math-macros}  % Custom package
\usepackage{math-theorems}  % Custom package
\usepackage{book-macros}  % Custom package for macros in the book

% Add subfiles --- each chapter can now be compiled by itself!
\usepackage{subfiles}

\makeatletter
\@ifclassloaded{EngC}{}{
    % Title page customization using titling package
    \usepackage{titling}
    \pretitle{
        \vspace{-2in}
    \begin{center}\LARGE\bfseries}
        \posttitle{
    \end{center}\vskip 0.5em}
    \preauthor{
    \begin{center}\large}
        \postauthor{
    \end{center}}
    \predate{
    \begin{center}\large}
        \postdate{
    \end{center}}

    % Customize vertical spacing
    \setlength{\droptitle}{-4em}

    % Add a rule under the title
    \renewcommand{\maketitlehooka}{%
        % \vspace{-2.75em}%
        \noindent\makebox[\linewidth]{\rule{\textwidth}{1pt}}%
        \vspace{3.5em}%
    }

    % Add space after author block
    \renewcommand{\maketitlehookb}{%
        \vspace{2em}%
    }

    % Add space before date
    \renewcommand{\maketitlehookc}{%
        \vspace{2em}%
    }

    % Add footnote content below date
    \renewcommand{\maketitlehookd}{%
        \vspace{6em}\footnotesize $^*$\,These authors contributed equally.
    }
}
\makeatother

\usepackage{etoolbox}
% Patches the command names after other packages are loaded
\patchcmd{\theorem}{Theorem}{Teorema}{}{}
\patchcmd{\example}{Example}{Ejemplo}{}{}

\begin{document}

\makeatletter
\@ifclassloaded{EngC}{
    \title{\Huge Principios y práctica del aprendizaje profundo de
    representaciones \\ {{\LARGE o Una teoría matemática de la memoria}}}
    \author{
        \textbf{Sam Buchanan}$^*$,
        Universidad de California, Berkeley \& TTIC  \\[\baselineskip]
        \textbf{Druv Pai}$^*$, Universidad de California, Berkeley
        \\[\baselineskip]
        \textbf{Peng Wang},
        Universidad de Macau \& Universidad de Michigan  \\[\baselineskip]
        \textbf{Yi Ma},
        Universidad de Hong Kong \& Universidad de California, Berkeley \\[3in]
        \footnotesize
        Este manuscrito se ha preparado para un curso sobre los principios del
        aprendizaje profundo
        y la inteligencia en general, a entregarse por los autores en
        sus respectivas
        instituciones. Esta versión previa a su publicación es de libre acceso para
        ver y descargar solo para
        uso personal y no se permiten su redistribución, su reventa ni su
        uso en obras
        derivadas.\\  Copyright \Pisymbol{psy}{227} 2026 Sam Buchanan,
        Druv Pai, Peng Wang, Yi
        Ma \\[2\baselineskip]
        $^*$\,Estos autores han contribuido en igual medida.
        %\footnotesize Copyright \Pisymbol{psy}{227}2020 Reserved \\ Ninguna
        %parte de este borrador se podrá reproducir o distribuir sin
        % autorización escrita
        %de los autores.\\
    }
}{
    \title{\huge Principios y práctica del \\
        \huge aprendizaje profundo de representaciones \\ \vspace{0.125in} \Large
        o Una teoría
    matemática de la memoria}

    \author{\vspace{0.5cm}
        \textbf{Sam Buchanan}$^*$,
        UC Berkeley \& TTIC  \\[\baselineskip]
        \textbf{Druv Pai}$^*$, UC Berkeley
        \\[\baselineskip]
        \textbf{Peng Wang},
        Universidad de Macau \& Universidad de Michigan  \\[\baselineskip]
        \textbf{Yi Ma},
        Universidad de Hong Kong \& UC Berkeley
        \vspace{5cm}\\
        \footnotesize
        Este manuscrito se ha preparado para un curso sobre los principios del
        aprendizaje profundo
        y la inteligencia en general, a entregarse por los autores en
        sus respectivas
        instituciones. Esta versión previa a su publicación es de libre acceso para
        ver y descargar solo para
        uso personal y no se permiten su redistribución, su reventa ni su
        uso en obras
        derivadas.\\  Copyright \Pisymbol{psy}{227} 2026 Sam Buchanan,
        Druv Pai, Peng Wang, Yi
        Ma
        %\footnotesize Copyright \Pisymbol{psy}{227}2020 Reserved \\ Ninguna
        %parte de este borrador se podrá reproducir o distribuir sin
        % autorización escrita
        %de los autores.\\
    }

    %Version 1.0: \date{\large 18 de enero de 2025}
    %Version 2.0:
    \date{\large 1° de marzo de 2026}
}
\makeatother

\frontmatter
\makeatletter
\@ifclassloaded{EngC}{}{
    \titlepage
    \thispagestyle{empty}
}
\makeatother
\maketitle

\phantomsection\addcontentsline{toc}{chapter}{\numberline{}Prefacio}
\subfile{chapters/preface/preface_es.tex}

\phantomsection\addcontentsline{toc}{chapter}{\numberline{}Prefacio
a la Versión 2.0}
\subfile{chapters/preface/preface-v2_es.tex}

\phantomsection\addcontentsline{toc}{chapter}{\numberline{}Declaración
de código abierto}
\subfile{chapters/declaration/declaration_es.tex}

\phantomsection\addcontentsline{toc}{chapter}{\numberline{}Agradecimientos}
\subfile{chapters/acknowledgment/acknowledgment_es.tex}

\phantomsection\addcontentsline{toc}{chapter}{\numberline{}Notación}
\subfile{chapters/notation_es.tex}

\renewcommand*{\contentsname}{Índice}
\tableofcontents

\setcounter{page}{1}
\pagenumbering{arabic}

\mainmatter

\subfile{chapters/chapter1/introduction_es}

\subfile{chapters/chapter2/classic-models_es}

\subfile{chapters/chapter3/denoising_es}

\subfile{chapters/chapter4/lossy-compression_es}

\subfile{chapters/chapter5/deep-networks_es}

\subfile{chapters/chapter6/consistent-representations_es}

\subfile{chapters/chapter7/conditional-inference_es}

\subfile{chapters/chapter8/applications_es}

\subfile{chapters/chapter9/future_es}

\makeatletter
\@ifclassloaded{EngC}{\backmatter}{\newpage}
\makeatother\appendix
\makeatletter
\@ifclassloaded{EngC}{}{
    \part*{Appendices}
    \addcontentsline{toc}{part}{Appendices}
}
\makeatother

\subfile{chapters/appendixA/optimization_es}

\subfile{chapters/appendixB/entropy_es}

\makeatletter
\@ifclassloaded{EngC}{\endappendix}{}
\makeatother

\addtocontents{toc}{\vspace{\baselineskip}}

\addcontentsline{toc}{chapter}{Bibliography}
\printbibliography

% \printindex{authors}{Author index}
% \printindex{subject}{Subject index}

\end{document}
